% Traced redraw of karabey_aksakalli_deployment_2021, Fig. 1 ("Conceptual
% model for microservices"): same boxes, same labels, same edges, same
% cardinality text -- a TikZ box-and-arrow rendering of a UML class
% diagram this project's own corpus already contains. This is the
% planted positive case bench_figure_similarity.py measures recall on.
% It carries no citekey, per every genre skill's figure rule.
\begin{tikzpicture}[
    box/.style={draw, rectangle, minimum width=3.2cm, minimum height=0.9cm, align=center, font=\small},
    >=Stealth
  ]
  \node[box] (app) at (0, 6) {MicroserviceApplication};
  \node[box] (infra) at (0, 3) {InfrastructureElement};
  \node[box] (inst) at (5, 3) {MicroserviceInstance};
  \node[box] (bare) at (-4.8, 0) {BareMetal};
  \node[box] (vm) at (-1.6, 0) {VirtualMachine};
  \node[box] (cont) at (1.6, 0) {Container};
  \node[box] (orch) at (4.8, 0) {OrchestrationPlatform};
  \node[box] (cross) at (9.5, 5.5) {CrossCuttingFunctionality};
  \node[box] (reg) at (9.5, 3) {ServiceRegistry};
  \node[box] (ext) at (9.5, 0.6) {ExternalSystem};
  \node[box] (comm) at (5, -3) {ServiceCommunicationEndPoint};

  \draw[<->] (app) -- node[right, font=\tiny] {[1..*] microserviceInstance} (inst);
  \draw[->] (app) -- node[left, font=\tiny] {[1..*] infrastructureElement} (infra);
  \draw[->] (inst) -- node[above, font=\tiny] {[0..*] deployedTo} (infra);
  \draw[->] (bare) -- (infra);
  \draw[->] (vm) -- (infra);
  \draw[->] (cont) -- (infra);
  \draw[->] (orch) -- (infra);
  \draw[->] (inst) -- node[above, font=\tiny] {[0..*] utilizes} (cross);
  \draw[->] (inst) -- node[right, font=\tiny] {[1..*] registeredTo} (reg);
  \draw[->] (inst) -- node[right, font=\tiny] {[0..*] integratesWith} (ext);
  \draw[->] (inst) -- node[below, font=\tiny] {[1..*] communicatesWith} (comm);
\end{tikzpicture}
